% Standalone-preprint adaptation of cusp_su2_magnetic.tex.
% Canonical source SHA-256: 8676F460516102827CCC32DAB613034F2153B8CE83EA87679AC6D6433121BE2A.
% The canonical radial-height symbol T was renamed R consistently; transpose
% notation was retained exactly.

\section{Integral magnetic sectors and a reducible
\texorpdfstring{\(SU(2)\)}{SU(2)} critical point on the radial cusp}
\label{cusp-magnetic-section}

Retain the positive radial flat tori
\(F_R=\mathbb R^4/P_R\mathbb Z^4\), their Euclidean quotient metrics, and
the real quantities \(L_R,q_R,m_R,D_R,H,B,m\) from
\cref{cusp-spec-radial-notation,cusp-spec-exact-dual-lattice,cusp-spec-D-expansion}.
In particular,
\[
 L_R=\operatorname{Im}\tau_R,\qquad
 q_R=-\operatorname{Im}\beta_R,\qquad
 D_R=L_Rq_R+6m_R^2=\operatorname{Vol}(F_R)>0.
\]
The real coordinate order is the already fixed order
\[
 x=(x_1,y_1,x_2,y_2)^{\mathsf t}.
\]

\subsection{The integral coframe}

Let \(e_1,\ldots,e_4\) be the standard basis of \(\mathbb R^4\), put
\(u=P_R^{-1}x\), and define
\begin{equation}
 \theta_R^j:=\mathrm du_j
 =(P_R^{-\mathsf T}e_j)^{\mathsf t}
   (\mathrm dx_1,\mathrm dy_1,\mathrm dx_2,\mathrm dy_2)^{\mathsf t}.
 \label{cusp-magnetic-integral-coframe}
\end{equation}

\begin{proposition}[Integral coframe and the distinguished unit-flux plane]
\label{cusp-magnetic-coframe-proposition}
The forms \(\theta_R^j\) are the integral coframe dual to the ordered
marked lattice basis:
\[
 \int_{[0,1]\ni s\mapsto[x+sP_Re_k]}\theta_R^j=\delta_{jk}.
\]
The first two members are exactly
\begin{equation}
 \theta_R^1=\frac{m_R}{D_R}\,\mathrm dy_1
             -\frac{L_R}{D_R}\,\mathrm dy_2,
 \qquad
 \theta_R^2=\frac{q_R}{D_R}\,\mathrm dy_1
             +\frac{6m_R}{D_R}\,\mathrm dy_2,
 \label{cusp-magnetic-first-two-covectors}
\end{equation}
and hence
\begin{equation}
 \theta_R^1\wedge\theta_R^2
 =\frac1{D_R}\,\mathrm dy_1\wedge\mathrm dy_2,
 \qquad
 \lvert\theta_R^1\wedge\theta_R^2\rvert^2
 =\frac1{D_R^2}.
 \label{cusp-magnetic-wedge-norm}
\end{equation}
\end{proposition}

\begin{proof}
For \(n\in\mathbb Z^4\),
\[
 u(x+P_Rn)=u(x)+n,
\]
which proves the integral-period assertion.  The first two columns of the
already-proved matrix \(P_R^{-\mathsf T}\) in
\cref{cusp-spec-full-inverse-transpose} give
\cref{cusp-magnetic-first-two-covectors}.  Their wedge coefficient is
\[
 \frac{6m_R^2+L_Rq_R}{D_R^2}=\frac1{D_R}.
\]
The coordinate forms \(\mathrm dy_1,\mathrm dy_2\) are orthonormal for the fixed quotient
metric, proving \cref{cusp-magnetic-wedge-norm}.
\end{proof}

\subsection{The unit-flux line and the topologically trivial rank-two sum}

The following construction is made in the marked coordinates \(u\) of
\eqref{cusp-magnetic-integral-coframe}; it therefore applies to the actual
skew lattice \(P_R\mathbb Z^4\), not to a rectangular replacement.

\begin{proposition}[Explicit factor of automorphy and connection]
\label{cusp-magnetic-line-bundle-proposition}
For \(n\in\mathbb Z^4\), set
\begin{equation}
 j_n(u)=\exp(-2\pi i n_1u_2).
 \label{cusp-magnetic-unit-factor}
\end{equation}
The quotient
\[
 \mathscr L_R=(\mathbb R^4\times\mathbb C)/\mathbb Z^4,
 \qquad
 n\cdot(u,z)=(u+n,j_n(u)z),
\]
is a Hermitian line bundle over \(F_R\).  The connection on its pullback
to \(\mathbb R^4\),
\begin{equation}
 \widetilde\nabla^{\mathscr L_R}
 =\mathrm d+2\pi i u_1\,\mathrm du_2,
 \label{cusp-magnetic-line-connection}
\end{equation}
descends, and its curvature is
\begin{equation}
 F_{\mathscr L_R}=2\pi i\,\mathrm du_1\wedge\mathrm du_2
        =2\pi i\,\theta_R^1\wedge\theta_R^2.
 \label{cusp-magnetic-line-curvature}
\end{equation}
With the unitary Chern--Weil convention
\(c_1=[\,iF/(2\pi)\,]\), one has
\(c_1(\mathscr L_R)=-[\theta_R^1\wedge\theta_R^2]\).
\end{proposition}

\begin{proof}
For \(m,n\in\mathbb Z^4\),
\[
 j_m(u+n)j_n(u)
 =j_{m+n}(u)\exp(-2\pi i m_1n_2)
 =j_{m+n}(u).
\]
Thus \eqref{cusp-magnetic-unit-factor} is a factor of automorphy.  Its
values are unitary, so the standard fibre metric descends.  If
\(\mathcal A=2\pi i u_1\,\mathrm du_2\), then
\[
 \mathcal A(u+n)=\mathcal A(u)+2\pi i n_1\,\mathrm du_2,
 \qquad
 j_n^{-1}\mathrm dj_n=-2\pi i n_1\,\mathrm du_2.
\]
Hence \(\mathcal A(u+n)+j_n^{-1}\mathrm dj_n=\mathcal A(u)\), which is precisely
the descent identity for \(\mathrm d+\mathcal A\) under the stated quotient
convention.  Taking \(\mathrm d\mathcal A\) gives
\eqref{cusp-magnetic-line-curvature}.  The Chern--Weil statement follows
without a sign suppression:
\[
 \frac{i}{2\pi}F_{\mathscr L_R}=-\theta_R^1\wedge\theta_R^2.
\]
\end{proof}

Let \(\mathscr L_R^{-1}=\mathscr L_R^*\) carry the dual metric and dual
connection, and set
\begin{equation}
 \mathscr E_R=\mathscr L_R\oplus\mathscr L_R^{-1}.
 \label{cusp-magnetic-rank-two-bundle}
\end{equation}

\begin{proposition}[Determinant, \(c_2\), and explicit
\(SU(2)\)-triviality]
\label{cusp-magnetic-su2-triviality}
The factors of automorphy of \(\mathscr E_R\) are
\[
 G_n(u)=\begin{pmatrix}j_n(u)&0\\0&j_n(u)^{-1}\end{pmatrix}\in SU(2).
\]
Evaluation gives a canonical parallel determinant trivialization
\[
 \det\mathscr E_R=\mathscr L_R\otimes\mathscr L_R^{-1}
 \xrightarrow{\ \mathrm{ev}\ }F_R\times\mathbb C.
\]
Moreover
\[
 c_1(\mathscr E_R)=0,
 \qquad
 c_2(\mathscr E_R)=0\quad\text{in }H^4(F_R;\mathbb Z),
\]
and the underlying principal \(SU(2)\)-bundle is smoothly trivial.  The
last assertion follows from the explicit nullhomotopy below and does not
rely only on the vanishing of a de Rham Chern form.
\end{proposition}

\begin{proof}
The determinant assertion follows pointwise from
\(j_nj_n^{-1}=1\).  The induced determinant connection has connection
form
\[
 \operatorname{tr}
 \begin{pmatrix}
  2\pi i u_1\,\mathrm du_2&0\\0&-2\pi i u_1\,\mathrm du_2
 \end{pmatrix}=0,
\]
so the evaluation trivialization is parallel.

Let \(c=c_1(\mathscr L_R)\).  The integral Whitney identity gives
\[
 c(\mathscr E_R)=(1+c)(1-c)=1-c^2.
\]
The factor \(j_n\) depends only on \((u_1,u_2)\), so \(\mathscr L_R\) is the
pullback of a line bundle \(L_0\) on
\(\mathbb R^2/\mathbb Z^2\) under the marked projection
\([u_1,u_2,u_3,u_4]\mapsto[u_1,u_2]\).  Hence
\[
 c^2=\pi_{12}^*\bigl(c_1(L_0)^2\bigr)=0
\]
integrally because \(H^4(\mathbb T^2;\mathbb Z)=0\).  This proves the two
Chern-class identities on the four-torus.

It remains to trivialize the \(SU(2)\)-bundle itself.  On the fundamental
square \(0\le u_1,u_2\le1\) for the first two marked coordinates, the only
nontrivial seam is
\[
 (0,t,v)\sim(1,t,\gamma(t)v),
 \qquad
 \gamma(t)=
 \begin{pmatrix}e^{-2\pi it}&0\\0&e^{2\pi it}\end{pmatrix}.
\]
For \(0\le r\le1\), define
\[
 a(r,t)=(1-r)e^{-2\pi it}+r,
 \qquad
 s(r)=\sin(\pi r),
 \qquad
 N(r,t)=\bigl(|a(r,t)|^2+s(r)^2\bigr)^{1/2},
\]
and
\begin{equation}
 U(r,t)=\frac1{N(r,t)}
 \begin{pmatrix}
  a(r,t)&-s(r)\\
  s(r)&\overline{a(r,t)}
 \end{pmatrix}.
 \label{cusp-magnetic-explicit-nullhomotopy}
\end{equation}
If \(0<r<1\), then \(s(r)>0\); at \(r=0,1\), one has
\(|a(r,t)|=1\).  Thus \(N\) never vanishes.  The two columns in
\eqref{cusp-magnetic-explicit-nullhomotopy} are orthonormal and its
determinant is one, so \(U(r,t)\in SU(2)\).  It is smooth and periodic in
\(t\), and
\[
 U(0,t)=\gamma(t),
 \qquad
 U(1,t)=I_2.
\]
This is an explicit nullhomotopy of the clutching loop.  To see directly
that it produces a smooth global frame, choose a smooth
\(\chi:[0,1]\to[0,1]\) which is zero near \(0\) and one near \(1\), and put
\[
 K(\sigma,t)=U(1-\chi(\sigma),t).
\]
Then \(K=I_2\) near \(\sigma=0\), \(K=\gamma\) near \(\sigma=1\), and
\(K(\sigma,0)=K(\sigma,1)\).  The required trivialization, written in the
direction from the clutched bundle to the product bundle, is
\begin{equation}
 [\sigma,t,v]\longmapsto
 \bigl([\sigma,t],K(\sigma,t)^{-1}v\bigr).
 \label{cusp-magnetic-explicit-global-frame}
\end{equation}
Indeed, the quotient convention gives
\([0,t,v]=[1,t,\gamma(t)v]\), and the two representatives have the same
image because
\[
 K(1,t)^{-1}\gamma(t)v
 =\gamma(t)^{-1}\gamma(t)v
 =v=K(0,t)^{-1}v.
\]
At the other seam,
\(K(\sigma,0)^{-1}=K(\sigma,1)^{-1}\).  The constant collars used in the
definition of \(\chi\) make the local matrix in
\cref{cusp-magnetic-explicit-global-frame} constant after applying the
seam transition, so the map and its inverse are smooth across the seam.
It is fibrewise unitary and determinant one.  Extending it independently
of \(u_3,u_4\) and
pulling it through the marking gives the required global
\(SU(2)\)-frame of \(\mathscr E_R\).
\end{proof}

\subsection{A non-flat reducible Yang--Mills critical connection}

The symbol \(H=\operatorname{Im}h(0)\) is already fixed by
\cref{cusp-spec-limiting-imaginary-parts}.  For the requested Lie-algebra
element, fix instead the unambiguous notation
\begin{equation}
 \langle X,Y\rangle_{\mathfrak{su}(2)}
   =-\frac12\operatorname{Tr}(XY),
 \qquad
 H_{\mathrm{mag}}=\begin{pmatrix}i&0\\0&-i\end{pmatrix}.
 \label{cusp-magnetic-su2-normalization}
\end{equation}
\begin{theorem}[Reducible critical point and exact action]
\label{cusp-magnetic-ym-critical-point}
The direct-sum connection on \(\mathscr E_R\) has, in the lifted splitting,
\begin{equation}
 \nabla^R=\mathrm d+2\pi u_1\,\mathrm du_2\,H_{\mathrm{mag}},
 \qquad
 F_{\nabla^R}
 =2\pi H_{\mathrm{mag}}\,\theta_R^1\wedge\theta_R^2
 =\frac{2\pi}{D_R}H_{\mathrm{mag}}\,\mathrm dy_1\wedge\mathrm dy_2.
 \label{cusp-magnetic-su2-curvature}
\end{equation}
After any of the global \(SU(2)\)-trivializations from
 \cref{cusp-magnetic-su2-triviality}, this is a connection on the
trivial principal \(SU(2)\)-bundle.  It is non-flat and satisfies
\[
 \mathrm d_{\nabla^R}^*F_{\nabla^R}=0.
\]
For the normalization
\begin{equation}
 \operatorname{YM}_R(\nabla)
 =\frac1{2g_{\mathrm{YM}}^2}
  \int_{F_R}\lvert F_\nabla\rvert^2\,
  \mathrm d\operatorname{vol}_{F_R},
 \label{cusp-magnetic-action-normalization}
\end{equation}
its action is exactly
\begin{equation}
 \operatorname{YM}_R(\nabla^R)
 =\frac{2\pi^2}{g_{\mathrm{YM}}^2D_R}.
 \label{cusp-magnetic-exact-action}
\end{equation}
Its first inverse-power correction, in the canonical scalar notation
\(H=\operatorname{Im}h(0)\), \(B=\operatorname{Im}b(0)\), is
\begin{equation}
 \operatorname{YM}_R(\nabla^R)
 =\frac{2\pi^2}{g_{\mathrm{YM}}^2R^2}
  \left(1+\frac{B-2H}{R}
  +O(R^{-2})\right).
 \label{cusp-magnetic-action-first-correction}
\end{equation}
\end{theorem}

\begin{proof}
The connection matrix is
\[
 \begin{pmatrix}
  2\pi i u_1\,\mathrm du_2&0\\
  0&-2\pi i u_1\,\mathrm du_2
 \end{pmatrix}
 =2\pi u_1\,\mathrm du_2\,H_{\mathrm{mag}},
\]
so \cref{cusp-magnetic-su2-curvature} follows from
\cref{cusp-magnetic-line-bundle-proposition,cusp-magnetic-wedge-norm}.
The diagonal endomorphism \(H_{\mathrm{mag}}\) is a
globally defined parallel section of \(\operatorname{ad}\mathscr E_R\),
because every transition \(G_n\) commutes with it.  The forms
\(\theta_R^j\) are parallel constant forms for the descended flat metric.
Thus \(\nabla^R F_{\nabla^R}=0\), which implies
\(\mathrm d_{\nabla^R}^*F_{\nabla^R}=0\).  The curvature is nonzero because
\(D_R>0\).

The chosen Lie-algebra normalization gives
\[
 \lvert H_{\mathrm{mag}}\rvert^2
 =-\frac12\operatorname{Tr}(H_{\mathrm{mag}}^2)=1.
\]
Equations \cref{cusp-magnetic-wedge-norm,cusp-spec-exact-volume} therefore
give
\[
 \int_{F_R}\lvert F_{\nabla^R}\rvert^2\,
   \mathrm d\operatorname{vol}_{F_R}
 =4\pi^2D_R\frac1{D_R^2}=\frac{4\pi^2}{D_R},
\]
which proves \cref{cusp-magnetic-exact-action}.  Finally,
\cref{cusp-spec-D-expansion} and the elementary reciprocal expansion
\[
 \frac1{R^2+aT+O(1)}
 =\frac1{R^2}\left(1-\frac aT+O(R^{-2})\right)
\]
with \(a=2H-B\) prove
\cref{cusp-magnetic-action-first-correction}.
\end{proof}

\subsection{Constant-curvature magnetic Laplacians for arbitrary integral
nondegenerate flux}

We next give the Hilbert space, closed form, self-adjoint realization, and
oscillator spectrum in a form valid for every nondegenerate integral
alternating form on the marked four-torus.

Let
\[
 \mathsf E=\sum_{1\le j<k\le4}E_{jk}\,\mathrm du_j\wedge\mathrm du_k
 \in\bigwedge^2(\mathbb Z^4)^*,
 \qquad E_{jk}\in\mathbb Z,
\]
and let the same symbol denote its skew matrix.  For \(n\in\mathbb Z^4\),
put
\begin{equation}
 J_n^{\mathsf E}(u)
 =\exp\!\left(-2\pi i\sum_{j<k}E_{jk}n_ju_k\right).
 \label{cusp-magnetic-general-factor}
\end{equation}
The integer in
\[
 \sum_{j<k}E_{jk}m_jn_k
\]
shows directly that
\(J_m^{\mathsf E}(u+n)J_n^{\mathsf E}(u)=J_{m+n}^{\mathsf E}(u)\).
Thus \eqref{cusp-magnetic-general-factor} defines a Hermitian line bundle
\(\mathscr L_{\mathsf E,R}\).  On its pullback, set
\begin{equation}
 \nabla^{\mathsf E}
 =\mathrm d+2\pi i\sum_{j<k}E_{jk}u_j\,\mathrm du_k.
 \label{cusp-magnetic-general-connection}
\end{equation}
Exactly as in \cref{cusp-magnetic-line-bundle-proposition}, the
connection descends and
\begin{equation}
 F_{\nabla^{\mathsf E}}=2\pi i\,\mathsf E.
 \label{cusp-magnetic-general-curvature}
\end{equation}

In \(u\)-coordinates the flat metric matrix is
\begin{equation}
 G_R=P_R^{\mathsf T}P_R,
 \qquad
 \sqrt{\det G_R}=D_R.
 \label{cusp-magnetic-flat-Gram}
\end{equation}
Define the Hilbert space concretely by
\begin{align}
 \mathcal H_{\mathsf E,R}
 =\{\,&f\in L^2_{\mathrm{loc}}(\mathbb R^4):
 f(u+n)=J_n^{\mathsf E}(u)f(u)\text{ for every }n\in\mathbb Z^4\},
 \label{cusp-magnetic-Hilbert-space}\\
 \lVert f\rVert^2
 ={}&\int_{[0,1]^4}\lvert f(u)\rvert^2\sqrt{\det G_R}\,\mathrm d^4u.
 \nonumber
\end{align}
Let \(\mathcal C_{\mathsf E,R}\) be the smooth functions with the same
automorphy.  On this dense subspace define
\begin{equation}
 q_{\mathsf E,R}[f]
 =\int_{[0,1]^4}
  (G_R^{-1})^{jk}
  \overline{\nabla^{\mathsf E}_jf}
  \nabla^{\mathsf E}_kf
  \sqrt{\det G_R}\,\mathrm d^4u.
 \label{cusp-magnetic-quadratic-form}
\end{equation}

\begin{theorem}[Friedrichs realization and oscillator spectrum]
\label{cusp-magnetic-oscillator-theorem}
The form \eqref{cusp-magnetic-quadratic-form} is nonnegative and closable.
Its closed form domain is
\[
 \operatorname{Dom}\overline q_{\mathsf E,R}
 =H^1(F_R;\mathscr L_{\mathsf E,R}).
\]
Its Friedrichs operator is the
magnetic Laplacian
\[
 \Delta_{\mathsf E,R}=(\nabla^{\mathsf E})^*\nabla^{\mathsf E},
 \qquad
 \operatorname{Dom}(\Delta_{\mathsf E,R})
 =H^2(F_R;\mathscr L_{\mathsf E,R}),
\]
and it has compact resolvent.

Assume now that \(\mathsf E\) is nondegenerate.  Let \(b_1,b_2>0\) be
defined by the eigenvalues
\begin{equation}
 \operatorname{Spec}\!\left(
  G_R^{-1/2}(2\pi\mathsf E)G_R^{-1/2}\right)
 =\{ib_1,-ib_1,ib_2,-ib_2\}.
 \label{cusp-magnetic-frequency-definition}
\end{equation}
Then, as a set (with finite multiplicities),
\begin{equation}
 \operatorname{Spec}(\Delta_{\mathsf E,R})
 =\left\{(2n_1+1)b_1+(2n_2+1)b_2:
 n_1,n_2\in\mathbb Z_{\ge0}\right\}.
 \label{cusp-magnetic-oscillator-spectrum}
\end{equation}
In particular, the ground-to-first-distinct-level spacing is
\begin{equation}
 \operatorname{gap}(\Delta_{\mathsf E,R})=2\min(b_1,b_2).
 \label{cusp-magnetic-level-gap}
\end{equation}
No multiplicity formula is needed for, or asserted in, these identities.
\end{theorem}

\begin{proof}
In local unitary trivializations, the integrand in
\eqref{cusp-magnetic-quadratic-form} is the squared norm of a first-order
unitary covariant derivative.  A finite trivializing cover and a partition
of unity show that the form norm
\(\lVert f\rVert^2+q_{\mathsf E,R}[f]\) is equivalent to the usual local \(H^1\)
norm.  Its completion is therefore
\(H^1(F_R;\mathscr L_{\mathsf E,R})\), proving closability and identifying
the closed domain.  The representation theorem for closed nonnegative
forms gives the Friedrichs operator.  Its local principal symbol is
\(G_R^{-1}(\xi,\xi)\), so it is uniformly elliptic.  Standard interior
elliptic regularity on the compact manifold without boundary identifies
the operator domain with \(H^2\).  The compact inclusion
\(H^1\hookrightarrow L^2\) implies compact resolvent.

The real skew-normal-form theorem supplies a constant orthonormal frame
\(v_1,\ldots,v_4\) on the Euclidean cover in which the real two-form
\(2\pi\mathsf E\) is
\[
 b_1v^1\wedge v^2+b_2v^3\wedge v^4.
\]
Changing the sign of \(v^2\) or \(v^4\), if necessary, makes both
frequencies positive.  Let \(\nabla_a\) denote covariant differentiation
along \(v_a\).  Each \(\nabla_a\) is formally skew-adjoint on the common
invariant core \(\mathcal C_{\mathsf E,R}\); the closures used below are the
ones determined by the closed magnetic form.  On that smooth core one has
\[
 [\nabla_1,\nabla_2]=ib_1,
 \qquad
 [\nabla_3,\nabla_4]=ib_2,
\]
and all other commutators vanish.  Put
\[
 C_1=\nabla_1-i\nabla_2,
 \qquad
 C_2=\nabla_3-i\nabla_4,
 \qquad
 \Delta_j=-(\nabla_{2j-1}^2+\nabla_{2j}^2).
\]
Direct multiplication, retaining the commutator signs, gives
\begin{equation}
 C_j^*C_j=\Delta_j-b_j,
 \qquad
 C_jC_j^*=\Delta_j+b_j,
 \qquad
 [\Delta_{\mathsf E,R},C_j]=-2b_jC_j.
 \label{cusp-magnetic-ladder-identities}
\end{equation}

Because the resolvent is compact, the lowest eigenvalue has a smooth
eigenfunction \(\psi_0\neq0\).  If \(C_j\psi_0\neq0\), the last identity
in \eqref{cusp-magnetic-ladder-identities} would make it an eigenfunction
with a strictly smaller eigenvalue.  Hence \(C_1\psi_0=C_2\psi_0=0\), and
the first identity gives
\[
 \Delta_{\mathsf E,R}\psi_0=(b_1+b_2)\psi_0.
\]
Conversely, repeated application of \(C_1^*,C_2^*\) produces nonzero
eigenfunctions with eigenvalues
\((2n_1+1)b_1+(2n_2+1)b_2\).  Nonvanishing follows from
\(C_jC_j^*=C_j^*C_j+2b_j\).

It remains to exclude other eigenvalues.  For an eigenfunction \(\psi\)
with eigenvalue \(\lambda>b_1+b_2\), the operator identity
\[
 (\lambda-b_1-b_2)\psi
 =C_1^*C_1\psi+C_2^*C_2\psi
\]
expresses \(\psi\) as a sum of raising operators applied to the lower
eigenfunctions \(C_j\psi\).  Induction over the finite set of eigenvalues
below \(\lambda\) reduces every eigenfunction to raised ground states.
This proves \eqref{cusp-magnetic-oscillator-spectrum}.  Raising once in a
least-frequency plane proves \eqref{cusp-magnetic-level-gap}.
\end{proof}

\begin{corollary}[Frequency product and cusp spacing bound]
\label{cusp-magnetic-frequency-product}
Let
\[
 d=|\operatorname{Pf}(\mathsf E)|
 =\lvert E_{12}E_{34}-E_{13}E_{24}+E_{14}E_{23}\rvert>0.
\]
Then
\begin{equation}
 b_1b_2=\frac{4\pi^2d}{D_R},
 \label{cusp-magnetic-product-identity}
\end{equation}
and
\begin{align}
 \operatorname{gap}(\Delta_{\mathsf E,R})
 &\le 4\pi\sqrt{\frac d{D_R}}
 \label{cusp-magnetic-gap-bound}\\
 &=\frac{4\pi\sqrt d}{R}
 \left(1+\frac{B-2H}{2R}
 +O(R^{-2})\right).
 \nonumber
\end{align}
\end{corollary}

\begin{proof}
The determinant of the skew matrix in
\eqref{cusp-magnetic-frequency-definition} is both
\(b_1^2b_2^2\) and
\[
 \frac{\det(2\pi\mathsf E)}{\det G_R}
 =\frac{(2\pi)^4\operatorname{Pf}(\mathsf E)^2}{D_R^2}.
\]
Taking the positive square root proves
\eqref{cusp-magnetic-product-identity}.  Since
\(\min(b_1,b_2)\le\sqrt{b_1b_2}\),
\Cref{cusp-magnetic-oscillator-theorem} gives
\eqref{cusp-magnetic-gap-bound}.  Expanding the square root of
\cref{cusp-spec-D-expansion} gives the displayed first correction.
\end{proof}

\subsection{A background-gauge negative mode}

We finish with the classical background-gauge Jacobi operator at the
reducible critical connection.  Every sign below uses exactly the inner
product and the element \(H_{\mathrm{mag}}\) fixed in
\cref{cusp-magnetic-su2-normalization}.

Put
\[
 \mathscr H_R^{\mathrm{ad},1}
 :=L^2\Omega^1(F_R;\operatorname{ad}\mathscr E_R).
\]
Put also
\[
 \mathscr H_R^{\mathrm{ad},0}
 :=L^2\Omega^0(F_R;\operatorname{ad}\mathscr E_R),
 \qquad
 D_R^{\max}:=\mathrm d_{\nabla^R}^{*,\max}.
\]
For the resolvent argument these real Hilbert spaces are complexified;
all operators commute with the real structure, so their restrictions to
real sections are unchanged.
Here the maximal distributional adjoint is the closed, densely defined
operator whose domain and value are specified without an implicit
regularity assumption by
\begin{align}
 \operatorname{Dom}(D_R^{\max})
 :=\{a\in\mathscr H_R^{\mathrm{ad},1}:{}&\text{ there is a }u\in
 \mathscr H_R^{\mathrm{ad},0}\text{ such that}\notag\\[-1mm]
 &\langle a,\mathrm d_{\nabla^R}\phi\rangle_{L^2}
 =\langle u,\phi\rangle_{L^2}
 \text{ for every }\phi\in
 C^\infty(F_R;\operatorname{ad}\mathscr E_R)\},
 \label{cusp-magnetic-maximal-codifferential-domain}\\
 D_R^{\max}a:={}&u.\notag
\end{align}
The test space is dense, so \(u\) is unique.  The definition also proves
closedness directly: if \(a_n\to a\) and \(D_R^{\max}a_n\to u\) in
\(L^2\), passage to the limit in the displayed testing identity gives
\(a\in\operatorname{Dom}(D_R^{\max})\) and \(D_R^{\max}a=u\).
Consequently the background slice
\begin{equation}
 \mathscr S_R
 :=\ker D_R^{\max}
 \subset\mathscr H_R^{\mathrm{ad},1}.
 \label{cusp-magnetic-closed-background-slice}
\end{equation}
is a closed subspace.

The scalar covariant Laplacian used below is fixed as follows.  On
\(C^\infty(F_R;\operatorname{ad}\mathscr E_R)\) start with the
nonnegative form
\[
 q_R^{(0)}[\phi,\psi]
 :=\langle\mathrm d_{\nabla^R}\phi,
             \mathrm d_{\nabla^R}\psi\rangle_{L^2}.
\]
Its closure has domain
\(H^1\Omega^0(F_R;\operatorname{ad}\mathscr E_R)\), and its associated
Friedrichs operator is
\begin{equation}
 \Delta_{\nabla^R}^{(0)}
 :=\mathrm d_{\nabla^R}^*\mathrm d_{\nabla^R},
 \qquad
 \operatorname{Dom}(\Delta_{\nabla^R}^{(0)})
 =H^2\Omega^0(F_R;\operatorname{ad}\mathscr E_R).
 \label{cusp-magnetic-scalar-covariant-laplacian-domain}
\end{equation}
The \(H^2\) domain follows from elliptic regularity on the compact flat
torus; in particular, this notation always denotes the nonnegative
self-adjoint realization, not merely its differential expression.

For an adjoint-valued one-form \(a\), define the smooth bundle
endomorphism
\[
 (V_Ra)_\nu:=-2\sum_{\mu=1}^4[F_{\nu\mu},a_\mu].
\]
It is symmetric: after interchanging \(\mu,\nu\), the identities
\(F_{\mu\nu}=-F_{\nu\mu}\) and skew-adjointness of
\(\operatorname{ad}(F_{\nu\mu})\) give
\(\langle V_Ra,b\rangle=\langle a,V_Rb\rangle\).
On the full one-form Hilbert space define
\begin{equation}
 \mathfrak j_R^{\mathrm{gf}}[a,b]
 :=\langle\nabla^Ra,\nabla^Rb\rangle_{L^2}
   +\langle V_Ra,b\rangle_{L^2},
 \qquad
 \operatorname{Dom}(\mathfrak j_R^{\mathrm{gf}})
 =H^1\Omega^1(F_R;\operatorname{ad}\mathscr E_R).
 \label{cusp-magnetic-full-jacobi-form}
\end{equation}
Because \(V_R\) is a bounded symmetric endomorphism on the compact
torus, this form is closed and bounded below.  Its associated
self-adjoint operator \(\mathcal J_R^{\mathrm{gf}}\) has
\begin{equation}
 \operatorname{Dom}(\mathcal J_R^{\mathrm{gf}})
 =H^2\Omega^1(F_R;\operatorname{ad}\mathscr E_R)
 \label{cusp-magnetic-full-jacobi-domain}
\end{equation}
and differential expression
\[
 (\mathcal J_R^{\mathrm{gf}}a)_\nu
 =-\sum_{\mu=1}^4\nabla_{e_\mu}\nabla_{e_\mu}a_\nu
  -2\sum_{\mu=1}^4[F_{\nu\mu},a_\mu].
\]

\begin{proposition}[Intertwining and the reducing Coulomb slice]
\label{cusp-magnetic-jacobi-reduction-proposition}
At the Yang--Mills connection \(\nabla^R\), the following formal
identities hold on smooth sections:
\begin{equation}
 \mathcal J_R^{\mathrm{gf}}\mathrm d_{\nabla^R}
 =\mathrm d_{\nabla^R}\Delta_{\nabla^R}^{(0)},
 \qquad
 \mathrm d_{\nabla^R}^*\mathcal J_R^{\mathrm{gf}}
 =\Delta_{\nabla^R}^{(0)}\mathrm d_{\nabla^R}^*.
 \label{cusp-magnetic-jacobi-intertwining}
\end{equation}
For every \(z\in\mathbb C\setminus\mathbb R\), the correctly typed
closed-operator resolvent identity is
\begin{equation}
 D_R^{\max}(\mathcal J_R^{\mathrm{gf}}-z)^{-1}f
 =(\Delta_{\nabla^R}^{(0)}-z)^{-1}D_R^{\max}f,
 \qquad f\in\operatorname{Dom}(D_R^{\max}).
 \label{cusp-magnetic-jacobi-resolvent-intertwining}
\end{equation}
In particular, \(\mathscr S_R\) reduces
\(\mathcal J_R^{\mathrm{gf}}\).  The part of the full operator in this
subspace and its closed form are therefore
\begin{equation}
 \mathcal J_R
 :=\mathcal J_R^{\mathrm{gf}}\big|_{\mathscr S_R},
 \quad
 \operatorname{Dom}(\mathcal J_R)
 =H^2\Omega^1(F_R;\operatorname{ad}\mathscr E_R)\cap\mathscr S_R,
 \quad
 \operatorname{Dom}(\mathfrak j_R)
 =H^1\Omega^1(F_R;\operatorname{ad}\mathscr E_R)\cap\mathscr S_R
 \label{cusp-magnetic-reduced-jacobi-domain}
\end{equation}
with
\(\mathfrak j_R=\mathfrak j_R^{\mathrm{gf}}|_{\operatorname{Dom}(\mathfrak j_R)}\).
The operator \(\mathcal J_R\) is self-adjoint on \(\mathscr S_R\), and
\(\mathfrak j_R\) is its closed, lower-bounded form.
\end{proposition}

\begin{proof}
Write \(\nabla_\mu=\nabla^R_{e_\mu}\) in the constant orthonormal frame,
so
\([\nabla_\mu,\nabla_\nu]\phi=[F_{\mu\nu},\phi]\).
For a smooth adjoint-valued scalar \(\phi\), commuting the derivatives
without suppressing either curvature contribution gives
\begin{align*}
 &\bigl(\mathcal J_R^{\mathrm{gf}}\mathrm d_{\nabla^R}\phi
       -\mathrm d_{\nabla^R}\Delta_{\nabla^R}^{(0)}\phi\bigr)_\nu\\
 &\quad=-\nabla_\mu\nabla_\mu\nabla_\nu\phi
         +\nabla_\nu\nabla_\mu\nabla_\mu\phi
         -2[F_{\nu\mu},\nabla_\mu\phi]\\
 &\quad=-[\nabla_\mu F_{\mu\nu},\phi]
   =[ (\mathrm d_{\nabla^R}^*F_{\nabla^R})_\nu,\phi].
\end{align*}
The last member vanishes by
\cref{cusp-magnetic-ym-critical-point}.  This proves the first identity
in \eqref{cusp-magnetic-jacobi-intertwining}.  Since both Laplacians and
\(\mathcal J_R^{\mathrm{gf}}\) are formally symmetric, pairing it against
arbitrary smooth zero- and one-forms proves the second identity.

Fix \(z\notin\mathbb R\), take
\(f\in\operatorname{Dom}(D_R^{\max})\), and put
\(u=(\mathcal J_R^{\mathrm{gf}}-z)^{-1}f\).
The full operator domain gives \(u\in H^2\), hence
\(D_R^{\max}u\in H^1\).  Moreover
\((\mathcal J_R^{\mathrm{gf}}-z)u=f\) and the distributional extension of
the second formal identity imply
\begin{equation}
 (\Delta_{\nabla^R}^{(0)}-z)D_R^{\max}u=D_R^{\max}f
 \quad\text{in }\mathcal D'(F_R;\operatorname{ad}\mathscr E_R).
 \label{cusp-magnetic-resolvent-distributional-equation}
\end{equation}
The right side is in \(L^2\).  Elliptic regularity for the scalar
Laplacian upgrades \(D_R^{\max}u\) from \(H^1\) to \(H^2\), so it belongs
to \(\operatorname{Dom}(\Delta_{\nabla^R}^{(0)})\).  Applying the bounded
scalar resolvent to
\eqref{cusp-magnetic-resolvent-distributional-equation} proves
\eqref{cusp-magnetic-jacobi-resolvent-intertwining}; no extension of
\(D_R^{\max}\) to all of \(L^2\) is asserted or needed.

If \(f\in\mathscr S_R\), the right side of
\eqref{cusp-magnetic-jacobi-resolvent-intertwining} is zero, hence both
resolvents at \(z\) and at \(\overline z\) preserve \(\mathscr S_R\).
For \(g\in\mathscr S_R^\perp\) and \(s\in\mathscr S_R\), self-adjointness
of the full operator gives
\[
 \langle(\mathcal J_R^{\mathrm{gf}}-z)^{-1}g,s\rangle
 =\langle g,(\mathcal J_R^{\mathrm{gf}}-\overline z)^{-1}s\rangle=0.
\]
Thus the resolvent preserves both \(\mathscr S_R\) and its orthogonal
complement, which is precisely the reducing property.  The standard
part of a self-adjoint operator on a reducing subspace is self-adjoint,
and its domain is the intersection in
\eqref{cusp-magnetic-reduced-jacobi-domain}.  The reducing orthogonal
projection commutes with the spectral calculus of the lower-bounded full
operator, including its form square root; this proves the asserted form
domain and completes the proof.
\end{proof}

\begin{theorem}[Exact charged Landau negative mode]
\label{cusp-magnetic-negative-mode}
For the self-adjoint Yang--Mills Jacobi operator \(\mathcal J_R\) on the
closed background slice \(\mathscr S_R\) defined above, the differential
expression in a constant orthonormal frame \((e_\mu)_{\mu=1}^4\) is
\begin{equation}
 (\mathcal J_Ra)_\nu
 =-\sum_{\mu=1}^4\nabla_{e_\mu}\nabla_{e_\mu}a_\nu
  -2\sum_{\mu=1}^4[F_{\nu\mu},a_\mu].
 \label{cusp-magnetic-jacobi-operator}
\end{equation}
There exists a nonzero smooth real
\(\mathfrak{su}(2)\)-valued one-form \(a\) such that
\begin{equation}
 \lVert a\rVert_{L^2}=1,
 \qquad
 \mathrm d_{\nabla^R}^*a=0,
 \qquad
 \mathcal J_Ra=-\frac{4\pi}{D_R}a.
 \label{cusp-magnetic-negative-eigen-equation}
\end{equation}
Equivalently, for the action normalization
\cref{cusp-magnetic-action-normalization},
\begin{equation}
 \operatorname{YM}_R''\big|_{\nabla^R}(a,a)
 =-\frac{4\pi}{g_{\mathrm{YM}}^2D_R}.
 \label{cusp-magnetic-negative-hessian}
\end{equation}
Thus the reducible non-flat critical connection has a strictly negative
classical second-variation direction for every finite \(R\).
\end{theorem}

\begin{proof}
First isolate the exact two-dimensional magnetic problem.  Let
\[
 W=\{x_1=x_2=0\}\subset\mathbb R^4,
\]
so that \((y_1,y_2)\) are oriented orthonormal coordinates on \(W\).
By \cref{cusp-spec-real-period-matrix}, orthogonal projection to \(W\)
sends the first two marked lattice generators to
\[
 v_1=(6m_R,-q_R),
 \qquad
 v_2=(L_R,m_R),
\]
and sends the last two to zero.  Since
\(\det(v_1,v_2)=D_R>0\), projection descends to
\[
 \varpi_R:F_R\longrightarrow
 \Sigma_R:=W/(\mathbb Zv_1+\mathbb Zv_2).
\]
The factors and connection of \(\mathscr L_R^2\) depend only on
\((u_1,u_2)\).  They therefore define a Hermitian line bundle with unitary
connection \((\mathscr M_R,\nabla^{\mathscr M_R})\) on \(\Sigma_R\) such
that
\begin{equation}
 (\mathscr L_R^2,\nabla^{\mathscr L_R^2})
 =\varpi_R^*(\mathscr M_R,\nabla^{\mathscr M_R}).
 \label{cusp-magnetic-descended-root-line}
\end{equation}
Let \(\Delta_+=(\nabla^{\mathscr M_R})^*\nabla^{\mathscr M_R}\) be the
nonnegative self-adjoint Bochner Laplacian on
\(L^2(\Sigma_R;\mathscr M_R)\), defined by its closed \(H^1\) form and
with operator domain \(H^2(\Sigma_R;\mathscr M_R)\).  More precisely, if
\(c_R>0\) is the constant Riemannian volume of a fibre of \(\varpi_R\),
then
\[
 U_R:L^2(\Sigma_R;\mathscr M_R)\longrightarrow
       L^2(F_R;\mathscr L_R^2),
 \qquad U_Rs=c_R^{-1/2}\varpi_R^*s,
\]
is an isometry.  Its image
\[
 \mathscr B_R:=U_RL^2(\Sigma_R;\mathscr M_R)
\]
is the closed fibre-basic subspace.  The pulled-back realization is,
by definition,
\[
 \Delta_{+,\mathrm{bas}}:=U_R\Delta_+U_R^{-1}
 \quad\text{on }\mathscr B_R,
 \qquad
 \operatorname{Dom}(\Delta_{+,\mathrm{bas}})
 =U_RH^2(\Sigma_R;\mathscr M_R),
\]
and its closed form domain is \(U_RH^1(\Sigma_R;\mathscr M_R)\).
Only this operator on \(\mathscr B_R\) is self-adjoint.  In the formulas
below we abbreviate \(\Delta_{+,\mathrm{bas}}\) again to \(\Delta_+\);
no self-adjoint operator on the full
\(L^2(F_R;\mathscr L_R^2)\) is thereby defined or asserted.  On this
root line the curvature is
\begin{equation}
 F_{\mathscr L_R^2}
 =4\pi i\,\theta_R^1\wedge\theta_R^2
 =i\kappa_R\,\mathrm dy_1\wedge\mathrm dy_2,
 \qquad
 \kappa_R=\frac{4\pi}{D_R}>0.
 \label{cusp-magnetic-root-field-strength}
\end{equation}

We now construct an exact lowest Landau section on this actual projected
lattice.  In the marked coordinates put
\[
 \mathcal D_1:=\partial_{u_1},
 \qquad
 \mathcal D_2:=\partial_{u_2}+4\pi i u_1,
\]
and define
\begin{equation}
 \zeta_R:=\frac{q_R-6im_R}{m_R+iL_R},
 \qquad
 v_R:=-\operatorname{Im}\zeta_R
     =\frac{D_R}{m_R^2+L_R^2}>0.
 \label{cusp-magnetic-theta-parameter}
\end{equation}
Indeed,
\[
 \operatorname{Im}\bigl((q_R-6im_R)(m_R-iL_R)\bigr)
 =-(L_Rq_R+6m_R^2)=-D_R,
\]
which also checks the sign \(v_R>0\).
The series
\begin{equation}
 \Theta_R(u_1,u_2)
 =\sum_{\ell\in\mathbb Z}
 \exp\!\left(
  -2\pi i\zeta_R(\ell+u_1)^2+4\pi i\ell u_2
 \right)
 \label{cusp-magnetic-theta-series}
\end{equation}
converges normally with all derivatives, because the absolute value of
its \(\ell\)-th summand is
\(\exp(-2\pi v_R(\ell+u_1)^2)\).  Reindexing the series and using
\(\ell\in\mathbb Z\) give
\begin{equation}
 \Theta_R(u_1+1,u_2)=e^{-4\pi i u_2}\Theta_R(u_1,u_2),
 \qquad
 \Theta_R(u_1,u_2+1)=\Theta_R(u_1,u_2).
 \label{cusp-magnetic-theta-automorphy}
\end{equation}
Thus \(\Theta_R\) is a smooth section of \(\mathscr L_R^2\).  Termwise
differentiation gives
\begin{equation}
 (\mathcal D_1+\zeta_R\mathcal D_2)\Theta_R=0.
 \label{cusp-magnetic-theta-marked-equation}
\end{equation}
The exact period block on the \((y_1,y_2)\)-plane is
\[
 \begin{pmatrix}6m_R&L_R\\-q_R&m_R\end{pmatrix}.
\]
Its inverse transpose gives the physical covariant derivatives
\begin{equation}
 \nabla_{y_1}^+
  =\frac{m_R\mathcal D_1+q_R\mathcal D_2}{D_R},
 \qquad
 \nabla_{y_2}^+
  =\frac{-L_R\mathcal D_1+6m_R\mathcal D_2}{D_R}.
 \label{cusp-magnetic-physical-root-derivatives}
\end{equation}
Consequently
\begin{equation}
 \nabla_{y_1}^+-i\nabla_{y_2}^+
 =\frac{m_R+iL_R}{D_R}
   (\mathcal D_1+\zeta_R\mathcal D_2).
 \label{cusp-magnetic-annihilator-conversion}
\end{equation}
Orthogonality of the \(u_2\)-Fourier modes also gives the nonzero norm
\[
 \int_0^1\!\!\int_0^1\lvert\Theta_R\rvert^2
 \,\mathrm du_1\,\mathrm du_2
 =\sum_{\ell\in\mathbb Z}\int_0^1
   e^{-4\pi v_R(\ell+u_1)^2}\,\mathrm du_1
 =\frac1{2\sqrt{v_R}}.
\]
Set
\[
 \phi_R:=\left(\frac{2\sqrt{v_R}}{D_R}\right)^{1/2}\Theta_R
\]
and pull it back by \(\varpi_R\).  Since the marked-coordinate volume
element is \(D_R\,\mathrm d^4u\), it has unit \(L^2(F_R)\)-norm.  Its
covariant derivatives in the \(x_1,x_2\) directions vanish, and
\begin{equation}
 (\nabla_{y_1}^+-i\nabla_{y_2}^+)\phi_R=0,
 \qquad
 \Delta_+\phi_R=\kappa_R\phi_R,
 \qquad
 \Delta_+\phi_R
 =-\bigl((\nabla_{y_1}^+)^2+(\nabla_{y_2}^+)^2\bigr)\phi_R.
 \label{cusp-magnetic-ground-equation}
\end{equation}
Indeed, \cref{cusp-magnetic-physical-root-derivatives} and
\([\mathcal D_1,\mathcal D_2]=4\pi i\) give
\[
 [\nabla_{y_1}^+,\nabla_{y_2}^+]
 =\frac{m_R(6m_R)-q_R(-L_R)}{D_R^2}\,4\pi i
 =i\kappa_R.
\]
Direct multiplication then gives
\[
 (\nabla_{y_1}^+-i\nabla_{y_2}^+)^*
 (\nabla_{y_1}^+-i\nabla_{y_2}^+)=\Delta_+-\kappa_R,
\]
and the second equation in \eqref{cusp-magnetic-ground-equation} follows
from the first.

In \(\mathfrak{sl}_2(\mathbb C)\), put
\[
 X_+=\begin{pmatrix}0&1\\0&0\end{pmatrix},
 \qquad
 X_-=\begin{pmatrix}0&0\\1&0\end{pmatrix}.
\]
Then
\[
 [H_{\mathrm{mag}},X_+]=2iX_+,
 \qquad
 [H_{\mathrm{mag}},X_-]=-2iX_-.
\]
The \(X_+\)-root line in
\(\operatorname{ad}(\mathscr L_R\oplus\mathscr L_R^{-1})
 \otimes\mathbb C\) is \(\mathscr L_R^2\),
so
\[
 a^+=\phi_RX_+(\mathrm dy_1-i\,\mathrm dy_2)
\]
is a globally defined complexified adjoint-valued one-form.  The first
equation in \eqref{cusp-magnetic-ground-equation} gives
\[
 \mathrm d_{\nabla^R}^*a^+
 =-(\nabla_{y_1}^+-i\nabla_{y_2}^+)\phi_R\,X_+=0.
\]

Put \(\mathcal B_R:=2\pi/D_R\), so that
\(\kappa_R=2\mathcal B_R\) and
\(F_{y_1y_2}=\mathcal B_RH_{\mathrm{mag}}\).  The two nonzero components
of \(a^+\) are \(a_{y_1}^+=\phi_RX_+\) and
\(a_{y_2}^+=-i\phi_RX_+\).  Substitution into
\eqref{cusp-magnetic-jacobi-operator}, with no suppressed charge factor,
gives
\begin{align*}
 (\mathcal J_Ra^+)_{y_1}
 &=\kappa_R\phi_RX_+
   -2[\mathcal B_RH_{\mathrm{mag}},-i\phi_RX_+]
  =-\kappa_R\phi_RX_+,\\
 (\mathcal J_Ra^+)_{y_2}
 &=-i\kappa_R\phi_RX_+
   -2[-\mathcal B_RH_{\mathrm{mag}},\phi_RX_+]
  =i\kappa_R\phi_RX_+.
\end{align*}
Thus \(\mathcal J_Ra^+=-\kappa_Ra^+\).

For completeness, formula \eqref{cusp-magnetic-jacobi-operator} follows by
expanding
\(F_{\nabla^R+ta}=F_{\nabla^R}+t\,\mathrm d_{\nabla^R}a+
\tfrac12t^2[a\wedge a]\), integrating the quadratic form by parts on
the background slice, and using the flat-base Weitzenbock identity.  The
Hodge Laplacian contributes
\(-[F_{\nu\mu},a_\mu]\), and the curvature term in the second variation
contributes the identical term, giving the coefficient \(-2\).

Finally take the anti-Hermitian real form
\begin{equation}
 a:=\frac1{\sqrt2}\bigl(a^+-(a^+)^\dagger\bigr)
  =\frac1{\sqrt2}\left(
   \phi_RX_+(\mathrm dy_1-i\,\mathrm dy_2)
   -\overline{\phi_R}X_-(\mathrm dy_1+i\,\mathrm dy_2)\right).
 \label{cusp-magnetic-real-negative-mode}
\end{equation}
It is nonzero, traceless, and anti-Hermitian.  For
\(X(z):=zX_+-\overline zX_-\), the fixed inner product gives
\(-\tfrac12\operatorname{Tr}(X(z)^2)=\lvert z\rvert^2\).  The two
orthonormal one-form components in
\cref{cusp-magnetic-real-negative-mode} therefore have total pointwise
norm \(\lvert a\rvert^2=\lvert\phi_R\rvert^2\), so
\(\lVert a\rVert_{L^2}=1\).  Complex conjugation and
Hermitian adjunction preserve both the gauge equation and the real Jacobi
operator, so \eqref{cusp-magnetic-negative-eigen-equation} follows from
the calculation for \(a^+\).  The Hessian identity follows from
\(\operatorname{YM}_R''(a,a)
=g_{\mathrm{YM}}^{-2}\langle a,\mathcal J_Ra\rangle_{L^2}\)
on the background slice.
\end{proof}


